\section{Extended frozen-policy ablations}
\label{app:ablations}

This appendix records the full design of experiment E5, the exact variant
specification strings passed to the ablation driver, and the per-opponent
breakdown that the summary in \S\ref{sec:ablations} compresses into a single
win rate. The grouped summary itself is Table~\ref{tab:ablations} and is not
repeated here. It also gives the E12 belief-substitution comparison under a
deliberately unfitted policy in full, and the method of the E16 declaration
pre-gate audit whose results are reported in \S\ref{sec:results-verify}.

\subsection{Design}
\label{app:abl-design}

The reference configuration is \vfast{} as shipped, the policy specification
\texttt{v04:mgate=0.008}, which equals the compiled defaults. Every variant
differs from it only in the switch named in its row; the \numParams{}-coordinate
fitted parameter vector is held at its shipping values throughout, and no
variant is refitted.

The experiment plays \numAblDeals{} deals against each member of a four-policy
panel --- \vthree{}, the turn-starvation lockout, the posterior detective and
\vtwo{} --- with each deal played in two seat rotations, giving \numAblGames{}
games per opponent per variant. The reported win rate is the mean over the four
opponents; the reference configuration scores \numAblRef\%. The seed bank is
606060, which is disjoint from every fitting and selection bank
(\S\ref{sec:fitting}). The command is

\begin{verbatim}
./fish ablate --ref=v04:mgate=0.008 --games=500 --rotations=2 --seed=606060 \
              --panel=v03,lockout,detective,v02 --variants="..."
\end{verbatim}

and the artifact is \filepath{research/v04/results/E5-ablations.json}.

Because the reference and every variant play the identical deals in the
identical rotations, the design is matched at the level of the deal. The
interval in the ``full $-$ ablated'' column is a paired percentile bootstrap in
which the \emph{deal}, not the game, is the resampling unit: deals are resampled
with replacement, the paired difference is recomputed on each resample, and the
endpoints are the 2.5\% and 97.5\% percentiles of \numBootstrapDraws{} draws.
Resampling
deals rather than games is the relevant correction here, because the six (or, in
E5, two) games generated from one deal are not independent observations.

\begin{center}
\fcolorbox{fishgreen}{fishlight}{%
\parbox{0.90\linewidth}{\vspace{3pt}
\textbf{Every row in E5 is a frozen-policy effect.} A row measures what happens
to \emph{this} fitted parameter vector when one component is removed or
replaced. It does not measure the standalone value of that component, and it
does not predict what a policy refitted without the component would score. No
matched-budget refitting experiment was run.
\vspace{3pt}}}
\end{center}

\subsection{Variant specification strings}
\label{app:abl-specs}

Table~\ref{tab:ablspecs} gives the specification string for each row exactly as
it is passed to \texttt{fish ablate}, so that any row can be reconstructed with
a single \texttt{fish match} invocation. The grammar is documented in
\S\ref{app:repro}. The rows of Table~\ref{tab:ablspecs} and of
Table~\ref{tab:ablopp} below carry the same ``Change'' labels, in the same
order, as Table~\ref{tab:ablations}, so a row can be followed across all three.

\begin{longtable}{@{}>{\footnotesize}p{0.42\linewidth}>{\footnotesize\ttfamily}p{0.52\linewidth}@{}}
\caption{E5 variant specifications. Each string is the complete policy
specification for one ablation row; the reference is
\texttt{v04:mgate=0.008}. \numAblDeals{} deals $\times$ 2 seat rotations
$=$ \numAblGames{} games against each of \vthree{}, the turn-starvation
lockout, the posterior detective and \vtwo{}; seed 606060; fitted weights
frozen throughout.}
\label{tab:ablspecs}\\
\toprule
\normalfont Change & \normalfont Specification string \\
\midrule
\endfirsthead
\toprule
\normalfont Change & \normalfont Specification string \\
\midrule
\endhead
\vfast{} (reference configuration) & v04:mgate=0.008 \\
\addlinespace
Drop capacity conditioning (C4) as well as the certificates & v04:mgate=0.008,belief=indep \\
Drop the ask-legality certificates (C5); plain Sinkhorn over C1--C4 & v04:mgate=0.008,belief=sinkhorn \\
Drop the ask-legality certificates (C5); exact capacity DP over C1--C4 & v04:mgate=0.008,belief=exact \\
\vblock{} belief substituted into the Fast-fitted policy & v04:mgate=0.008,belief=block \\
Fast belief with the policy prior off (matches Block's prior-free belief) & v04:mgate=0.008,belief=fast,ptheta=0,pphi=0 \\
Remove the fitted policy-aware prior & v04:mgate=0.008,ptheta=0,pphi=0 \\
\addlinespace
Zero the lock-completion weight & v04:mgate=0.008,w5=0 \\
Zero the hit-probability weight & v04:mgate=0.008,w0=0 \\
Remove the two-ply refinement & v04:mgate=0.008,topk=0 \\
Zero both information-leak weights & v04:mgate=0.008,w9=0,w19=0 \\
Remove the one-ply value lookahead & v04:mgate=0.008,value=0 \\
Zero the reply-threat weight & v04:mgate=0.008,w8=0 \\
Zero the runway weight & v04:mgate=0.008,w18=0 \\
\addlinespace
Declaration threshold 0.99 & v04:mgate=0.008,decl=0.99 \\
Name the allocation by conditioned greedy MAP & v04:mgate=0.008,gmap=1 \\
Fixed threshold in place of the value-based stopping rule & v04:mgate=0.008,vdecl=0 \\
Declaration threshold 0.80 & v04:mgate=0.008,decl=0.80 \\
Cash locked half-suits immediately & v04:mgate=0.008,patient=0 \\
\bottomrule
\end{longtable}

The \texttt{belief=fast,ptheta=0,pphi=0} row exists so that the belief-substitution
row has a prior-matched comparator: the reference engine carries no fitted
policy prior, so switching to it changes two things at once, and this row
isolates the prior half of the change. Because \texttt{belief=fast} is already
the default, it is numerically identical to the row that sets
\texttt{ptheta=0,pphi=0} alone. The \texttt{patient=0} row is inert by
construction: when the
value-based stopping rule is active it decides first, so the immediate-cash
switch is never consulted, and the paired difference is exactly zero with a
zero-width interval. Both are reported for completeness rather than as
evidence.

\subsection{Belief-mode variants are not all exact}
\label{app:abl-beliefmodes}

Three rows sit between \numAblCFourAbs\% and \numAblCFiveAbs\% and it would be
easy to read them as ablations of exact inference. They are not.
\texttt{belief=exact} is the card-level capacity
dynamic program \emph{without} the ask-legality certificates (C5);
\texttt{belief=sinkhorn} is plain Sinkhorn scaling with neither certificates nor
the disjointness handling; \texttt{belief=indep} additionally drops capacity
conditioning (C4), leaving independent per-card ownership estimates. Their
paired differences from the reference are \numAblCFive{}
(\numAblCFiveCI{}), \numAblSinkhorn{} (\numAblSinkhornCI{}) and \numAblCFour{}
(\numAblCFourCI{}) points respectively, each a paired percentile bootstrap over
deals. The exact
reference engine of \S\ref{sec:belief} is \texttt{belief=block} alone, and it
appears in the row at \numAblBlockAbs\%. The name \texttt{exact} in the specification
grammar is a legacy identifier for the capacity dynamic program and is
unfortunate; it is retained because it appears in the committed artifact.

\subsection{Per-opponent breakdown}
\label{app:abl-peropponent}

Table~\ref{tab:ablopp} gives the four per-opponent win rates behind each
averaged row, read directly from
\filepath{research/v04/results/E5-ablations.json}. The artifact records
per-opponent win rates but not per-opponent intervals, so no uncertainty is
attached to the individual columns; only the averaged difference carries the
paired interval reported in Table~\ref{tab:ablations}. With \numAblGames{}
games behind each cell the resolution of a cell is 0.1 percentage points. The
averaged win rate of each row is not repeated here, because it is the ``win
rate'' column of Table~\ref{tab:ablations}; the paired difference is repeated so
that each row's overall effect sits beside its per-opponent split, and the
averaged win rate is \numAblRef\% minus that difference.

{\footnotesize
\begin{longtable}{@{}p{0.42\linewidth}rrrrr@{}}
\caption{E5 per-opponent breakdown. Win rate of the frozen \vfast{}
configuration or of the named variant against each panel member;
\numAblDeals{} deals $\times$ 2 seat rotations $=$ \numAblGames{} games per
cell, seed 606060, fitted weights frozen in every row. Rows carry the same
labels and order as Table~\ref{tab:ablations}. ``Full $-$ ablated'' is the
averaged paired difference of that table, so the row's overall win rate is
\numAblRef\% minus it. Per-opponent cells carry no interval; the artifact
reports intervals only for the averaged paired difference. Source:
\texttt{research/v04/results/E5-ablations.json}.}
\label{tab:ablopp}\\
\toprule
Change & Full $-$ ablated & \vthree{} & Lockout & Detective & \vtwo{} \\
\midrule
\endfirsthead
\toprule
Change & Full $-$ ablated & \vthree{} & Lockout & Detective & \vtwo{} \\
\midrule
\endhead
\vfast{} (reference configuration) & --- & 75.0\% & 78.3\% & 72.8\% & 83.9\% \\
\midrule
Drop capacity conditioning (C4) as well as the certificates & \numAblCFour{} & 21.1\% & 20.9\% & 19.0\% & 24.9\% \\
Drop the ask-legality certificates (C5); plain Sinkhorn over C1--C4 & \numAblSinkhorn{} & 20.0\% & 23.2\% & 24.1\% & 31.6\% \\
Drop the ask-legality certificates (C5); exact capacity DP over C1--C4 & \numAblCFive{} & 27.2\% & 26.5\% & 28.1\% & 32.6\% \\
\vblock{} belief substituted into the Fast-fitted policy & \numAblBlock{} & 69.2\% & 70.4\% & 69.6\% & 76.0\% \\
Fast belief with the policy prior off (matches Block's prior-free belief) & \numAblPrior{} & 69.8\% & 75.7\% & 69.6\% & 79.4\% \\
Remove the fitted policy-aware prior & \numAblPrior{} & 69.8\% & 75.7\% & 69.6\% & 79.4\% \\
\midrule
Zero the lock-completion weight & \numAblLockWeight{} & 71.3\% & 74.1\% & 69.6\% & 82.2\% \\
Zero the hit-probability weight & \numAblHitWeight{} & 76.9\% & 78.4\% & 69.2\% & 79.3\% \\
Remove the two-ply refinement & \numAblTwoPly{} & 73.3\% & 77.2\% & 73.2\% & 81.4\% \\
Zero both information-leak weights & \numAblLeak{} & 74.9\% & 76.7\% & 74.7\% & 83.2\% \\
Remove the one-ply value lookahead & \numAblValue{} & 77.3\% & 79.5\% & 72.4\% & 82.1\% \\
Zero the reply-threat weight & \numAblThreat{} & 74.7\% & 78.9\% & 74.8\% & 83.0\% \\
Zero the runway weight & \numAblRunway{} & 75.1\% & 76.7\% & 76.0\% & 83.6\% \\
\midrule
Declaration threshold 0.99 & \numAblDeclNinetyNine{} & 74.4\% & 78.5\% & 72.8\% & 83.5\% \\
Name the allocation by conditioned greedy MAP & \numAblGmap{} & 74.6\% & 78.2\% & 73.1\% & 83.4\% \\
Fixed threshold in place of the value-based stopping rule & \numAblStop{} & 76.1\% & 78.0\% & 73.5\% & 81.9\% \\
Declaration threshold 0.80 & \numAblDeclEighty{} & 75.0\% & 78.0\% & 73.1\% & 83.8\% \\
Cash locked half-suits immediately & $+0.00$ & 75.0\% & 78.3\% & 72.8\% & 83.9\% \\
\bottomrule
\end{longtable}
}

The three degraded-belief rows are degraded against all four opponents
together, by between forty and sixty percentage points in every cell of the
table above; no single panel member drives the averaged effect. Several of the
small-effect rows are not uniform
in sign across the panel. Zeroing the hit-probability weight raises the win rate
against \vthree{} (75.0\% to 76.9\%) and against the lockout (78.3\% to 78.4\%)
while lowering it against the detective (72.8\% to 69.2\%) and \vtwo{} (83.9\%
to 79.3\%); zeroing the runway weight raises it against the detective (72.8\% to
76.0\%) while lowering it against the lockout (78.3\% to 76.7\%). Where the
averaged interval already includes zero, these sign reversals are consistent
with sampling variation, and the per-opponent cells are not powered to
distinguish that from a real opponent-specific effect.

\subsection{Belief substitution under a deliberately unfitted policy (E12)}
\label{app:abl-e12}

The belief-substitution row of Table~\ref{tab:ablations} compares two beliefs
under weights that were fitted around one of them. E12 repeats the substitution
under a policy chosen so that neither belief has been fitted to: pure
posterior-greedy asking, with the single ask weight $w_0 = 12$ and every other
ask weight zero, no value term, no two-ply refinement, and no policy prior. The
two specifications are

\begin{verbatim}
v04:weights=12|0|0|0|0|0|0|0|0|0|0|0|0|0|0|0|0|0|0|0,\
    value=0,topk=0,ptheta=0,pphi=0,belief=fast

v04:weights=12|0|0|0|0|0|0|0|0|0|0|0|0|0|0|0|0|0|0|0,\
    value=0,topk=0,ptheta=0,pphi=0,belief=block
\end{verbatim}

(written on one line each in \filepath{engine/experiments.sh}; the backslash above
is a line break for the page only). Each plays \numETwelveDeals{} deals in six
seat rotations, \numETwelveGames{} games, against \vthree{} on seed bank 959595,
under the full dialect.

The Fast belief scores \numETwelveFast\% and the reference belief
\numETwelveBlock\%. The 95\% intervals are \numETwelveFastCI\% and
\numETwelveBlockCI\%, each a single-arm deal-clustered percentile bootstrap. The
two arms share the seed bank, so the deals are matched, but the artifact reports
those single-arm intervals rather than an interval on the difference; the paired
quantity was not computed for this experiment.

The Fast arm's point estimate exceeds the Block arm's by just over one
percentage point, in the same direction as
the fitted comparison of Table~\ref{tab:ablations}. The intervals overlap
substantially, so E12 does not resolve the question, and it does not establish
that the direction of the fitted result is independent of the ask weights having
been fitted around the Fast path. The experiment that would resolve it ---
refitting the entire policy against the reference belief under a matched
optimisation budget and comparing the two fitted agents --- was not run, and is
listed as a limitation in \S\ref{sec:limitations}.

\subsection{Method of the declaration pre-gate audit (E16)}
\label{app:abl-gateaudit}

\begin{table}[!htbp]
\centering
\caption{Declaration pre-gate false-negative audit (E16). Population: the primary
bank of \numHtHDeals{} held-out deals in six seat rotations, seed $90210$, all
eight panel opponents, \vfast{} with policy weights frozen; the unit is one
declaration opportunity, or for the first three rows one (opportunity, live
half-suit) pair. Shares are exact fractions of the stated denominator; no
intervals are constructed.}
\label{tab:gateaudit}
\small
\begin{tabular}{lrr}
\toprule
Quantity & Count & Share \\
\midrule
Declaration opportunities examined & 10{,}102{,}149 & --- \\
(opportunity, live half-suit) pairs & 58{,}530{,}669 & --- \\
Pairs rejected by a cheap gate & 24{,}114{,}786 & 41.20\% of pairs \\
\midrule
False negatives (rejected, would have been declared) & 1{,}017 & 0.00422\% of rejections \\
Declarations made, gated path & 332{,}645 & --- \\
Declarations made, ungated path & 333{,}661 & --- \\
Opportunities where the chosen action differs & 1{,}016 & 0.0101\% of opportunities \\
\bottomrule
\end{tabular}
\end{table}


The results of E16 are in \S\ref{sec:results-verify} and
Table~\ref{tab:gateaudit}; this subsection records how the audit is constructed.

\vfast{} screens declaration candidates with two cheap scores before the
expensive posterior query of \S\ref{sec:declare} runs, and the two screens are
separate because they cut at different stages and read different information.
The first is \texttt{Knowledge::cheapTeamProb} in
\filepath{engine/src/belief.hpp}, a capacity-only product: for each unresolved
card of the half-suit it takes the share of the surviving support's total
capacity that belongs to the declarer's team, and multiplies. It uses no dynamic
program and no Sinkhorn iteration, so it can be evaluated before the posterior
is refreshed at all. \texttt{proposeDeclaration} in
\filepath{engine/src/v04.hpp} applies it twice: once across all live half-suits,
to decide whether to refresh the posterior for this opportunity at all, and once
per half-suit, against the threshold \texttt{gateTeamProb}. The second screen
lives inside \texttt{evaluateSet} and is a product of per-card team
probabilities read off the refreshed posterior marginals, tested against
\texttt{marginalGate}. It therefore prunes a half-suit that survived the
capacity-only screen but whose posterior team mass is small, before the joint
allocation query is run. Both thresholds are $0.008$ in the shipping
configuration, and both are bypassed when the unresolved pool has fallen to
eight cards or fewer, or when the forced-endgame pressure level is non-zero.

With \texttt{gateaudit=1} the policy carries out a second, ungated evaluation
alongside the one it acts on. At every declaration opportunity, and for
\emph{every} live half-suit rather than only the ones that survive the screens,
\texttt{evaluateSet} is re-run with both gates disabled and the result is passed
to the same stopping rule \eqref{eq:stop} at the same pressure level. The
audited path changes no action: the declaration the policy actually makes is
still the gated one, so the games played are the games of the primary bank.

A (opportunity, half-suit) pair counts as \emph{rejected} when the gated path
did not reach a full verdict on it --- because no half-suit passed the
first screen at that opportunity, or because this half-suit failed the
per-half-suit screen, or because the gated \texttt{evaluateSet} returned no
verdict while the ungated one did. A rejected pair counts as a \emph{false
negative} when the ungated evaluation plus the same stopping rule would have
declared it. Separately, an opportunity is counted as one where the chosen
action differs when the gated and ungated paths disagree about whether to
declare at all, or agree to declare but select different half-suits. The two
counts are not the same quantity: several pairs can be falsely rejected at one
opportunity, and a falsely rejected pair changes nothing if the gated path
declared a different half-suit at that opportunity anyway. Both counts, and the
number of declarations each path made, are in Table~\ref{tab:gateaudit}.

The audit is run by

\begin{verbatim}
PANEL=v03,lockout,detective,v02,diversifier,hunter,bluffer,random
./fish gateaudit --games=700 --rotations=6 --seed=90210 --panel=$PANEL
\end{verbatim}

which replays the primary bank of \numHtHDeals{} deals in six seat rotations
against all eight panel opponents with the fitted weights frozen, and
accumulates the counters across every seat and every game. Every figure it
reports is an exact count over that corpus; no interval is constructed, and the
measurement is specific to this policy, this bank and the two thresholds
compiled into the shipping configuration.
